% Generated draft: requires mathematical and source review.
\documentclass[11pt]{article}
\usepackage[margin=25mm]{geometry}
\usepackage{fontspec}
\setmainfont{Latin Modern Roman}
\usepackage{amsmath,amssymb,mathtools,unicode-math}
\setmathfont{Latin Modern Math}
\usepackage{enumitem,needspace,xurl,hyperref}
\setlength{\parindent}{0pt}
\setlength{\parskip}{5pt}
\setlength{\emergencystretch}{3em}
\newcommand{\N}{\mathbb{N}}
\newcommand{\Z}{\mathbb{Z}}
\newcommand{\Q}{\mathbb{Q}}
\newcommand{\R}{\mathbb{R}}
\newcommand{\C}{\mathbb{C}}
\newcommand{\F}{\mathbb{F}}
\newcommand{\A}{\mathbb{A}}
\newcommand{\PP}{\mathbb{P}}
\newcommand{\E}{\mathbb{E}}
\newcommand{\eps}{\varepsilon}
\newcommand{\dd}{\,\mathrm{d}}
\begin{document}
\section*{UnsolvedMath ID: 701212}
\textbf{Source problem code:} catalogue entry
\subsection{Definitions and mathematical statement}
A quasivariety of groups is a class axiomatized, within groups, by quasi-identities: universally quantified implications whose premises are finite conjunctions of equations and whose conclusions are equations. Let \(M\) be such a class, let \(A\in M\), and let \(H\leq A\). The source defines
\[
\operatorname{dom}^{M}_{A}(H)=\{a\in A:\ \forall B\in M\ \forall f,g\in\operatorname{Hom}(A,B),\ f|_H=g|_H\Longrightarrow f(a)=g(a)\}.
\]
There is an unresolved domain qualification when this notation is subsequently applied to every quasivariety \(N\subseteq M\): the source does not require \(A\in N\), although its initial definition requires membership in the ambient quasivariety. To display a precise candidate formulation, define, for every quasivariety \(N\subseteq M\),
\[
D_N(A,H)=\{a\in A:\ \forall B\in N\ \forall f,g\in\operatorname{Hom}(A,B),\ f|_H=g|_H\Longrightarrow f(a)=g(a)\}.
\]
This formula makes sense even when \(A\notin N\); using it in that case is an explicit extension of the source's stated domain convention, requiring review.

Under this extension, put
\[
\mathcal L(M,A,H)=\{D_N(A,H):N\subseteq M\text{ is a quasivariety of groups}\},
\]
ordered by inclusion. A lattice here means that every pair has a greatest lower bound \(X\wedge Y\) and a least upper bound \(X\vee Y\) within \(\mathcal L(M,A,H)\). These operations are order-theoretic bounds in this family; they are not stipulated to be set intersection and union.

The question is whether there exist \(M\), \(A\in M\), and \(H\leq A\) such that \(\mathcal L(M,A,H)\) is a lattice, is modular, and is not distributive. Explicitly, modularity requires
\[
\forall X,Y,Z\in\mathcal L(M,A,H),\qquad
X\subseteq Z\Longrightarrow X\vee(Y\wedge Z)=(X\vee Y)\wedge Z,
\]
whereas failure of distributivity requires
\[
\exists X,Y,Z\in\mathcal L(M,A,H),\qquad
X\wedge(Y\vee Z)\ne (X\wedge Y)\vee(X\wedge Z).
\]
If the source instead intends dominions only when \(A\in N\), the family must be replaced by
\[
\mathcal L_A(M,A,H)=\{\operatorname{dom}^{N}_{A}(H):N\subseteq M\text{ is a quasivariety of groups and }A\in N\}.
\]
The supplied record does not resolve which family is intended. Neither formulation is asserted here to be equivalent to the other.
\subsection{Short English statement}
Can the inclusion lattice of dominions of a subgroup, obtained by varying the quasivariety inside a fixed quasivariety of groups, be modular but non-distributive? The source leaves unclear whether the varying quasivarieties must contain the given group. The supplied open status coexists with an explicitly unverified AI claim of an affirmative answer.
\subsection{Sources}
\begin{itemize}
\item Supplied local UnsolvedMath registry record, permanent ID \(701212\), exact source code \texttt{catalogue entry}, titled \emph{Kourovka Notebook Problem 16.20}; supplied upstream revision \texttt{d371ee917ac9647fbb06cd6fc529deec6f3dda82}.
\item A. I. Budkin, problem 16.20, \emph{The Kourovka Notebook (Unsolved Problems in Group Theory)}, 16th issue (2006). The record also cites the 21st issue edition (2026), identified as \texttt{arXiv:1401.0300v48}; its status notes specify October 2026. The supplied collection description attributes editorship to E. I. Khukhro and V. D. Mazurov. These are supplied bibliographic attributions, not independently checked citations.
\item The background and status notes report an explicitly labelled ``Unverified AI claim'' of an affirmative answer, attributed to M. Kinyon and J. Urban, \emph{Preprint}, 2026, at \url{https://kourovkanotebookorg.wordpress.com/wp-content/uploads/2026/09/kourovka-mathematics.pdf}. The status notes say the notebook marks this with \(\#\). No verified solution is supplied.
\item The supplied status review is dated \texttt{2026-10-16}, retains the canonical label \texttt{open}, and classifies its evidence as \texttt{unrefereed\_claim}. This assessment is supplied dated curation, not independently verified literature assessment. No source was visited for this draft.
\end{itemize}
\end{document}
